\section{Evaluation}\label{sec:eval}

This section evaluates the effectiveness of \tool at exposing consensus-critical SSZ canonicalization vulnerabilities.
We focus on two research questions:
\begin{itemize}
    \item \textbf{RQ1 (Reproduction):} Can \tool reliably reproduce historical SSZ canonicalization vulnerabilities on consensus-layer schemas?
    \item \textbf{RQ2 (Zero-Day):} Can \tool uncover previously unknown vulnerability instances in production SSZ implementations?
\end{itemize}

\subsection{Experimental Setup}

\paragraph{Benchmarks}
We compile a benchmark suite of \textbf{five} historical SSZ vulnerabilities, each originally discovered by \tool and later confirmed by maintainers (see our bug reports).
The suite spans five specification-defined canonicalization boundaries in SSZ~\cite{sszspec}:
(1) \emph{dirty bitvector padding} acceptance in \icode{Bitvector[N]};
(2) \emph{dirty boolean} encodings beyond \icode{0x00/0x01};
(3) \emph{null-bitlist} deserialization without the sentinel bit;
(4) \emph{union trailing-data} acceptance (notably relevant to \icode{Optional[T]} encodings in \cite{eip6800}); and
(5) \emph{variable-length container gaps} violating offset contiguity.

\paragraph{Targets, metric, and protocol}
Each benchmark is instantiated on an Ethereum consensus-layer schema that contains the affected field (e.g., \icode{BeaconState.JustificationBits} for dirty bitvector padding).
We report \emph{Time to Exposure (TTE)}, i.e., the wall-clock time required to produce the first failure-inducing input in a cold-start run.
Unless otherwise stated, we report the median TTE over 30 independent runs.

\paragraph{Baseline}
As a baseline, we use a structure-aware LibFuzzer~\cite{libfuzzer} configuration under a 30-minute budget per target.
Under this setting, LibFuzzer failed to reproduce any of the historical bugs.

\paragraph{Bug provenance and counting}
To avoid cherry-picking and to keep ground truth auditable, every benchmark and zero-day instance reported in this section corresponds to a historical bug recorded in \texttt{docs/bugs} (each initially discovered by \tool and subsequently confirmed upstream).
When counting zero-days, we count \emph{per-implementation} instances: a single specification boundary (e.g., dirty bitvector padding) may manifest as multiple independent bugs across different client libraries, and each vulnerable codebase is counted as a distinct instance.

\paragraph{Environment}
All measurements were collected on a commodity workstation (Apple M3 Max, macOS 15.6.1) using Go 1.25.1.

\subsection{RQ1: Reproducing Historical Vulnerabilities}

\autoref{tab:repro-report} summarizes reproduction results.
Across all five benchmarks, \tool exposes the vulnerability in a cold-start setting with millisecond-scale TTE.
In each case, \tool synthesizes a non-canonical but semantically valid encoding that triggers either a round-trip mismatch (i.e., \({\sf enc}({\sf dec}(w)) \neq w\)) or a hash-tree-root divergence under canonicalization.

\begin{table*}[t]
\centering
\caption{Bug reproduction effectiveness on historical SSZ vulnerabilities. We report median TTE over 30 cold-start runs. ``--'' indicates the bug was not found within the 30-minute timeout.}
\label{tab:repro-report}
\small
\resizebox{\textwidth}{!}{%
\begin{tabular}{lcccc}
\toprule
\textbf{Benchmark ID} & \textbf{Affected Implementations} & \textbf{Consensus Schema} & \textbf{LibFuzzer TTE (ms)} & \textbf{\tool TTE (ms)} \\
\midrule
BV-DirtyPadding & Prysm/fastssz, Teku, py-ssz, dynamic-ssz~{\scriptsize\cite{prysm,fastssz,teku,pyssz,dynamic_ssz}} & \icode{BeaconState.JustificationBits} & -- & 5.390 \\
Bool-Dirty      & Prysm/fastssz, dynamic-ssz~{\scriptsize\cite{prysm,fastssz,dynamic_ssz}} & \icode{Validator.Slashed} & -- & 0.351 \\
Bitlist-Null    & dynamic-ssz~{\scriptsize\cite{dynamic_ssz}} & \icode{PendingAttestation.AggregationBits} & -- & 4.020 \\
Union-DirtyTail & Lighthouse, Lodestar, remerkleable~{\scriptsize\cite{lighthouse,lodestar,remerkleable}} & \icode{SuffixStateDiff.new\_value} & -- & 0.164 \\
Gap-Offsets     & remerkleable~{\scriptsize\cite{remerkleable}} & \icode{PendingAttestation} & -- & 0.588 \\
\bottomrule
\end{tabular}
}
\end{table*}


These results indicate that \tool{}'s specification-aligned buckets expose canonicalization boundaries directly, enabling reliable reproduction even when the triggering inputs are semantically silent (i.e., no crash).

\subsection{RQ2: Zero-Day Instances in Production Implementations}

We further ran \tool on production-grade SSZ implementations used across major consensus clients and libraries.
\autoref{tab:0day-instances-report} lists the previously unknown vulnerability instances uncovered in this study.
Each row corresponds to a distinct vulnerable implementation (i.e., we count per-implementation instances, even when multiple projects share the same root cause).

\begin{table*}[t]
\centering
\caption{Previously unknown vulnerability instances discovered by \tool in production SSZ implementations.}
\label{tab:0day-instances-report}
\small
\begin{tabular}{cll}
\toprule
\textbf{\#} & \textbf{Vulnerability} & \textbf{Implementation / Issue} \\
\midrule
1 & Container deserialize short/empty input & Lodestar~\cite{lodestar} \href{https://github.com/ChainSafe/ssz/issues/502}{ChainSafe/ssz\#502} \\
2.1 & Bitvector dirty padding & Teku~\cite{teku} \href{https://github.com/Consensys/teku/pull/10173}{Consensys/teku\#10173} \\
2.2 & Bitvector dirty padding & Prysm (fastssz)~\cite{prysm,fastssz} \\
2.3 & Bitvector dirty padding & ethereum/py-ssz~\cite{pyssz} \\
2.4 & Bitvector dirty padding & dynamic-ssz~\cite{dynamic_ssz} \href{https://github.com/pk910/dynamic-ssz/issues/54}{pk910/dynamic-ssz\#54} \\
3.1 & Boolean canonicalization \{0x00,0x01\} & fastssz / Prysm~\cite{fastssz,prysm} \href{https://github.com/ferranbt/fastssz/issues/222}{ferranbt/fastssz\#222} \\
3.2 & Boolean canonicalization \{0x00,0x01\} & dynamic-ssz~\cite{dynamic_ssz} \href{https://github.com/pk910/dynamic-ssz/issues/53}{pk910/dynamic-ssz\#53} \\
4.1 & Union trailing data accepted & Lighthouse~\cite{lighthouse} \href{https://github.com/sigp/ethereum_ssz/issues/64}{sigp/ethereum\_ssz\#64} \\
4.2 & Union trailing data accepted & Lodestar~\cite{lodestar} \href{https://github.com/ChainSafe/ssz/issues/507}{ChainSafe/ssz\#507} \\
4.3 & Union trailing data accepted & ethereum/remerkleable~\cite{remerkleable} \href{https://github.com/protolambda/remerkleable/issues/28}{protolambda/remerkleable\#28} \\
5 & Variable-length container gap & ethereum/remerkleable~\cite{remerkleable} \href{https://github.com/protolambda/remerkleable/issues/29}{protolambda/remerkleable\#29} \\
6 & Null-bitlist trap & dynamic-ssz~\cite{dynamic_ssz} \href{https://github.com/pk910/dynamic-ssz/issues/56}{pk910/dynamic-ssz\#56} \\
\bottomrule
\end{tabular}
\end{table*}

In total, \tool uncovered \textbf{12} previously unknown vulnerability instances.
Several instances are consensus-breaking because they enable non-canonical encodings that are accepted by one client but rejected (or normalized) by another, yielding different state roots.

\subsection{Discussion}

These results highlight two practical implications.
First, consensus risk concentrates at a small set of specification-defined canonicalization boundaries, but each boundary can manifest as multiple independent implementation bugs.
Second, by systematically abstracting and prioritizing these boundaries, \tool provides a practical mechanism for continuously vetting heterogeneous SSZ stacks ahead of new spec features such as \cite{eip6800}.
